\documentclass[10pt,a4paper]{article}
\usepackage[margin=1.55cm]{geometry}
\usepackage{xcolor}
\usepackage{amsmath,amssymb}
\usepackage{graphicx}
\usepackage{enumitem}
\usepackage{fontspec}
\usepackage[CJKmath=true]{xeCJK}
% CJK main font is resolved at COMPILE time on the actual machine, not baked in at
% generation time, so a .tex generated on one OS still renders on another (the older
% Python-side probe hardcoded one OS's font name into the file and broke when moved).
% Walk a cross-platform priority list and use the first font that exists; AutoFakeBold
% + AutoFakeSlant give bold/italic CJK even on faces that ship no bold/italic cut.
\newcommand{\setCJKto}[1]{\setCJKmainfont{#1}[AutoFakeBold=true,AutoFakeSlant=0.2]}
\IfFontExistsTF{Noto Sans CJK SC}{\setCJKto{Noto Sans CJK SC}}{%        % Linux (fonts-noto-cjk); cross-platform if installed
 \IfFontExistsTF{Source Han Sans SC}{\setCJKto{Source Han Sans SC}}{%   % Adobe 思源黑体; cross-platform
  \IfFontExistsTF{PingFang SC}{\setCJKto{PingFang SC}}{%               % macOS 10.11+
   \IfFontExistsTF{Hiragino Sans GB}{\setCJKto{Hiragino Sans GB}}{%    % older macOS
    \IfFontExistsTF{Microsoft YaHei}{\setCJKto{Microsoft YaHei}}{%     % Windows
     \IfFontExistsTF{WenQuanYi Micro Hei}{\setCJKto{WenQuanYi Micro Hei}}{% % older Linux
      \setCJKto{Songti SC}}}}}}}                                       % macOS bottom fallback
\definecolor{cblue}{HTML}{4C72B0}
\setlength{\parskip}{3pt}\setlength{\parindent}{0pt}
\setlist[enumerate]{topsep=2pt,itemsep=2pt,parsep=1pt,partopsep=0pt}
\setlist[itemize]{topsep=2pt,itemsep=2pt,parsep=1pt,partopsep=0pt}
% Let prose-embedded inline math break across lines and absorb small overflows, so simple
% formulas inserted mid-sentence stop running past the right margin.
\tolerance=1200
\emergencystretch=2.5em
\relpenalty=20
\binoppenalty=40
\hfuzz=1pt
\begin{document}

\section*{EffectSlice: Eligibility-Gated Certificates for Task-Local Procedural Effects}
\textbf{Method.} EffectSlice
\section*{Motivation}
A system can reproduce a paper and still leave a user with a large procedure that is expensive or unnecessary for the user's own task. A short version is not automatically useful: the complete version may never have helped, the removed steps may matter only on the user's data, or the apparent gain may disappear after execution cost and safety regressions are counted.

PaperBench and AutoReproduce help evaluate or automate research reproduction. A Framework for Evaluating Agentic Skills at Scale measures whole-skill utility, and newer systems can generate, trace, retrieve, or compress skill procedures. What is still missing is a fair decision rule for one independently chosen task: first prove that the complete skill helps, then ask whether a strict nonempty subset keeps the same measured benefit.

EffectSlice makes that decision explicit. It turns procedures into deterministic source-linked units, keeps dependencies exact through strongly connected component (SCC) condensation, and tests candidate subsets with fresh paired runs. It issues a certificate only when the compact version stays within \(\epsilon \) of the complete version, preserves contracts and guardrails, and survives a corrected audit of every retained-unit deletion. The same gate is applied to outputs from other compact-skill methods. The result is not a claim that one reducer is universally best;\allowbreak{} it is a task-local answer to adopt, reduce, or abstain.
\section*{Method}
\subsection*{M1 Eligibility Gate}
\emph{Establish that the complete artifact has a sealed positive task-local net effect before reduction.}
\begin{enumerate}[leftmargin=*,start=1]
  \item Write one canonical JavaScript Object Notation (JSON) file before any run. Give every task case a stable identifier, immutable input and expected-output references, and a weight;\allowbreak{} identify the complete artifact F and same-scaffold baseline B;\allowbreak{} pin the versions of all scorers, contract checks, guardrail checks, parsers, and the executor;\allowbreak{} record the plain-language decision thresholds and error-allocation rules;\allowbreak{} and list separate random-seed identifiers for eligibility, discovery, and confirmation. Define one inference interface that converts paired case-and-seed rows into the full-effect estimate, its one-sided bound, the slice-to-full interval, and corrected pass/fail decisions. 【author decision: specify scorer direction and aggregation, every contract and guardrail predicate, missing-value behavior, the paired resampling unit, the simultaneous-bound method, and which decisions share each error family.】 Sort the JSON keys, normalize paths and numbers, hash the file and all referenced files with Secure Hash Algorithm 256-bit (SHA-256), and stop if any digest changes.
\par\emph{Why:} The decision rules cannot move after favorable or unfavorable slice results appear.
  \item For every sealed eligibility random seed and task case, create one pair identifier and run two isolated branches: the complete artifact F and the same-scaffold baseline B. Reset mutable state between branches and give both branches exactly the same case input and seed. Save the artifact digest, the raw and normalized scorer results, every contract result, every guardrail result, execution status, and an error-log digest. Match rows only by the pair identifier, reject missing or duplicate branch rows and unregistered retries, and produce an append-only ledger showing that every pair has one valid F row and one valid B row.
\par\emph{Why:} Reduction matters only relative to a measured full intervention on the user's task.
  \item Within each pair identifier, subtract the baseline B score from the complete-artifact F score in the direction declared by the scorer and keep the task-case weight. Send all paired differences to the frozen inference procedure. This produces the full-artifact task effect \(\Delta _\tau (F),\) its simultaneous one-sided lower confidence bound \(L_{F},\) and the preregistered minimum benefit \(\delta _{min}.\) Continue only when \(L_{F}\) is strictly above \(\delta _{min}\) and every corrected contract and guardrail check passes. Store the manifest and ledger hashes, the estimate and bound, the threshold, all predicate decisions, and the supporting row identifiers;\allowbreak{} otherwise return \(`\mathit{no\_validated\_beneficial\_effect}`\) and do not create slice candidates.
\noindent\emph{The complete artifact is eligible only when the simultaneous lower confidence bound of its cost-adjusted effect exceeds the preregistered minimum benefit.}
\par\noindent{\small\ttfamily [eq~1, raw LaTeX, not typeset] L\_F = \textbackslash\{\}operatorname\{LCB\}\_\{1-\textbackslash\{\}alpha\_F\}\textbackslash\{\}!\textbackslash\{\}left(\textbackslash\{\}Delta\_\textbackslash\{\}tau(F)\textbackslash\{\}right) > \textbackslash\{\}delta\_\{\textbackslash\{\}min\}}
\par\emph{Why:} This prevents an ineffective complete artifact from lending credibility to an ineffective subset.
\end{enumerate}
\subsection*{M2 Exact Procedural Object}
\emph{Compile auditable source units and reduce the feasible family through an exact dependency quotient.}
\begin{enumerate}[leftmargin=*,start=4]
  \item Parse the pinned SKILL.md by explicit structure only: heading-delimited leaf sections, list items, fenced executable blocks, and tool-template records. Describe each \(`source_{map}`\) span by file digest and exact byte and line intervals, and describe each executable region by pinned file digest, symbol or template identifier, and exact interval. Combine one workflow leaf, one uniquely intersecting source span, one uniquely intersecting executable region, and one contract role into an ordered atom. Give the atom a digest-derived identifier and store all four parts, source order, and component digests. Split only at explicit boundaries, never merge units because their meanings look similar, and return \(`ambiguous_{atomization}`\) with the competing matches if any part is missing or non-unique.
\par\emph{Why:} Every retained or removed procedure remains traceable and reproducible.
  \item Make one graph vertex for every frozen atom. Propose a dependency edge from procedure u to procedure v when u declares v as a prerequisite, reads an artifact produced by v, or the frozen executor reports that dependency while running the complete artifact. The edge \(u\to v\) means u cannot run correctly without v's output or uniquely attributable state change. Check the edge by replaying u from the same saved pre-state with v present and with only v and its attributable outputs hidden, while keeping every other prerequisite fixed. Keep the edge only if the second replay requests v, loses a required input, or fails a declared contract. Save the evidence and replay digests, combine duplicate evidence, and abstain if v cannot be isolated or an observed dependency has no represented producer.
\par\emph{Why:} A compact subset must still be executable, not just look semantically coherent.
  \item Run Tarjan's algorithm on dependency graph G to find Strongly Connected Components (SCCs), meaning groups whose procedures depend on one another in a cycle. Give each SCC a stable identifier from its sorted atom identifiers, replace each dependency between two groups by one edge, and call the resulting graph D. Map any dependency-closed atom subset to the SCCs it fully contains with \(\Phi ,\) and expand SCC identifiers back to atoms with Union. Since D may have no more than 16 nodes, enumerate every downward-closed subset and check that mapping to atoms and back preserves identifiers, dependency closure, the materialized artifact digest, and executor resolution. Recognize \(\Delta _{D}(T)\) and \(\Delta _{G}(Union(T))\) as the same measured effect only when both resolve to the identical atom set and artifact digest. Return all maps and checks, or \(`\mathit{granularity\_out\_of\_scope}`\) when D has more than 16 nodes.
\noindent\emph{The SCC projection and inverse form an exact correspondence between dependency-closed atom subsets and downward-closed quotient subsets.}
\par\noindent{\small\ttfamily [eq~2, raw LaTeX, not typeset] \textbackslash\{\}Phi(S)=\textbackslash\{\}\{C\textbackslash\{\}in\textbackslash\{\}operatorname\{SCC\}(G):C\textbackslash\{\}subseteq S\textbackslash\{\}\},\textbackslash\{\}quad \textbackslash\{\}operatorname\{Union\}(\textbackslash\{\}Phi(S))=S}
\noindent\emph{Quotienting changes neither the feasible artifact represented by a subset nor its measured task effect.}
\par\noindent{\small\ttfamily [eq~3, raw LaTeX, not typeset] \textbackslash\{\}Delta\_D(T)=\textbackslash\{\}Delta\_G(\textbackslash\{\}operatorname\{Union\}(T))}
\par\emph{Why:} The search becomes smaller without changing which executable artifacts can be certified.
  \item Let N be the number of SCC nodes in verified graph D and define Q as at most 24 registered candidate-and-neighbor tests. A candidate is valid only if its sorted SCC set is nonempty, leaves out at least one SCC from the complete artifact F, contains every prerequisite required by each retained SCC, expands to atoms with unchanged source records and digests, resolves in the executor, and passes every static contract check. Reserve the last N test indices only for the eventual deletion neighbors and allow discovery to use only the first \(Q- N\) indices. Give each index one append-only ledger row containing its purpose, hypothesis identifier, candidate digest, random-seed block, and status. Reject unregistered tests, reused indices, extra discovery tests, or any run that would consume an index reserved for a neighbor.
\par\emph{Why:} The final audit budget cannot be consumed by exploratory search.
\end{enumerate}
\subsection*{M3 Slice Admission}
\emph{Generate candidates independently of fresh outcomes and certify a locked locally irreducible slice.}
\begin{enumerate}[leftmargin=*,start=8]
  \item Call any frozen slice proposer through one adapter. Give it the candidate-registry digest, the already revealed discovery ledger, and its own frozen-state digest;\allowbreak{} require it to return one sorted SCC set \(T_{q}\) and, optionally, a trace. First register the query number, a fresh paired random-seed block, and a preassigned error share, but do not reveal the new \(T_{q},\) F, or B outcomes or raw scorer results. Check \(T_{q}\) with the candidate predicate from step 7 (Let N be the number of SCC nodes \(in\dots ),\) hash the proposer's name, version, settings, state, inputs, and output, and save the proposal before any branch runs. An invalid, duplicate, late, or outcome-informed proposal uses its registered test index and is recorded as rejected instead of being silently replaced.
\par\emph{Why:} Candidate generation stays separate from the admission test and its unseen evidence.
  \item For registered query q, run \(T_{q},\) F, and B on the same new pair identifiers with isolated state and the ledger format from step 2 (For every sealed eligibility \(random\dots ).\) This step first defines the paired slice-to-full effect gap \(G_\tau (T_{q}),\) its simultaneous confidence interval \(I_\tau (T_{q}),\) and the discovery margin \(M_\tau (T_{q})\);\allowbreak{} that margin is positive only when the whole interval is strictly inside the equivalence band whose half-width is \(\epsilon .\) Assign the query's error share \(\alpha _{D,q}\) by the preregistered predictable sequence and never move error shares between the full-artifact, discovery, and confirmation families after seeing results. Correct every contract and guardrail decision, using \(\eta _{j}\) as guardrail j's declared tolerance. Save the supporting pair rows, estimates, interval ends, margin, allocations, decisions, and correction details, and retain \(T_{q}\) only when \(M_\tau (T_{q})\) is strictly positive and all corrected checks pass.
\noindent\emph{The discovery margin is positive only when the entire simultaneous confidence interval for the slice-to-full effect gap lies inside the equivalence band.}
\par\noindent{\small\ttfamily [eq~4, raw LaTeX, not typeset] G\_\textbackslash\{\}tau(T)=\textbackslash\{\}Delta\_\textbackslash\{\}tau(\textbackslash\{\}operatorname\{Union\}(T))-\textbackslash\{\}Delta\_\textbackslash\{\}tau(F),\textbackslash\{\}quad M\_\textbackslash\{\}tau(T)=\textbackslash\{\}epsilon-\textbackslash\{\}sup\_\{g\textbackslash\{\}in I\_\textbackslash\{\}tau(T)\}|g|}
\noindent\emph{Predictable discovery allocations and sealed family budgets control repeated adaptive queries without reallocating error after outcomes.}
\par\noindent{\small\ttfamily [eq~5, raw LaTeX, not typeset] \textbackslash\{\}alpha\_\{D,q\}=\textbackslash\{\}frac\{6\textbackslash\{\}alpha\_D\}\{\textbackslash\{\}pi\textasciicircum{}2q\textasciicircum{}2\},\textbackslash\{\}quad \textbackslash\{\}alpha\_F+\textbackslash\{\}alpha\_D+\textbackslash\{\}alpha\_C\textbackslash\{\}leq\textbackslash\{\}alpha\_\{\textbackslash\{\}mathrm\{total\}\}}
\par\emph{Why:} Selection depends on the downstream effect gap to F, not an internal compactness proxy.
  \item Choose \(T_{star}\) from the surviving candidates with one frozen total ordering and save the exact discovery-ledger snapshot and ordering key used. 【author decision: state whether the ordering first prefers fewer retained SCCs, a larger discovery margin, or another criterion, and define a complete tie-break because the choice changes which neighbor family is called irreducible.】 For each SCC c kept in \(T_{star},\) form its deletion neighbor by removing c and every kept SCC that depends on c through any path;\allowbreak{} this leaves the largest dependency-closed subset that excludes c. Merge duplicate neighbor subsets but keep the list of original c identifiers. Test every neighbor on its reserved index with the corrected effect, contract, and guardrail rules. Every retained c must point to a tested neighbor that fails at least one rule;\allowbreak{} otherwise return an insufficiency abstention. Save \(T_{star},\) the selection evidence, the full origin-to-neighbor map, and all decisions.
\par\emph{Why:} The certificate cannot call a retained procedure necessary while leaving an unchecked deletion.
\end{enumerate}
\subsection*{M4 Sealed Comparison}
\emph{Apply one corrected gate to controls and baselines, then report held-out decisions including abstentions.}
\begin{enumerate}[leftmargin=*,start=11]
  \item Before reading untouched outcomes, create a sealed confirmation hypothesis ledger \(H_{C}\) with one row for F, B, \(T_{star},\) every registered deletion neighbor, the permuted-margin control, and each frozen direct output from ClawTrace, Graph-of-Skills, and SkillRAE. Send every compact output through the same atom compiler and dependency, source-validity, executor, and contract checks used for EffectSlice. 【author decision: define a method-neutral conversion for baseline outputs that lack exact source or executable-region records, and decide in advance whether an output that cannot be represented is rejected, counted as an abstention, or excluded by a stated applicability rule.】 Create the control by using the manifest's sealed permutation to shuffle candidate labels inside each paired discovery block without breaking branch pairs, and lock the resulting artifact before confirmation. Rerun every ledger row on untouched pair identifiers and use the frozen confirmation correction. Admit a compact output only when F is still eligible, the output is strict and nonempty, its slice-to-full interval is inside the equivalence band, every contract and guardrail passes, and every registered deletion neighbor is rejected. Return \(Cert_\tau \) with all identifying digests for a passing \(T_{star},\) or a reason-coded rejection or abstention for every ledger row.
\par\emph{Why:} A common admission protocol isolates the certificate from whichever reducer proposed the compact artifact.
  \item After development on exactly eight existing tasks, tag the source revision, hash every manifest and executable artifact, disable changes to settings, and accept only paper-task pair manifests whose identifiers and source snapshots did not appear during development. Register at least four new pairs before seeing any outcomes. 【author decision: define the eligible population of paper-task pairs, how pairs are selected or sampled, all exclusions, and an outcome-independent reference that decides whether an admitted slice is false;\allowbreak{} without these choices the held-out and false-admission rates cannot be reproduced.】 Run the unchanged step 1 (Write one canonical \(JavaScript\dots )–step 11\) (Before reading untouched outcomes) pipeline once for each pair and append one final row containing eligibility, atomization result, quotient size, candidate status, confirmation decision, abstention reason, and evidence digests. Calculate correctly admitted compact-artifact, false-admission, and abstention rates over every registered pair, with separate counts for ineligible, ambiguous, more-than-16-node, empty, and insufficient-test cases. Publish all pair-level reason codes and never remove a pair after its result is known.
\par\emph{Why:} Held-out pairs and visible abstentions test whether the decision is useful beyond the development cases.
\end{enumerate}
\end{document}
